% ============================================================
%  Seminar-Präsentation – Generisches Template
%  Beamer + Metropolis-Theme
% ============================================================
\documentclass[aspectratio=169, 12pt]{beamer}

% ── Theme ───────────────────────────────────────────────────
\usetheme{metropolis}
\usecolortheme{default}

% ── Sprache & Encoding ──────────────────────────────────────
\usepackage[T1]{fontenc}
\usepackage[utf8]{inputenc}
\usepackage[ngerman]{babel}   % ggf. auf [english] wechseln

% ── Schriften ───────────────────────────────────────────────
\usepackage{FiraSans}         % passt gut zu Metropolis
\usepackage{FiraMono}         % Monospace für Code

% ── Mathematik ──────────────────────────────────────────────
\usepackage{amsmath, amssymb, amsthm}

% ── Code-Highlighting ───────────────────────────────────────
\usepackage{listings}
\lstset{
  language     = SQL,
  basicstyle   = \ttfamily\small,
  keywordstyle = \bfseries\color{mLightBrown},
  commentstyle = \itshape\color{gray},
  stringstyle  = \color{mLightGreen},
  numbers      = left,
  numberstyle  = \tiny\color{gray},
  frame        = single,
  breaklines   = true,
  showstringspaces = false,
  tabsize      = 2,
}

% ── Grafik & Tabellen ───────────────────────────────────────
\usepackage{graphicx}
\usepackage{booktabs}
\usepackage{tikz}
\usetikzlibrary{positioning, arrows.meta}

% ── Bibliographie ───────────────────────────────────────────
\usepackage[backend=biber, style=acmnumeric]{biblatex}
\addbibresource{references.bib}   % eigene .bib-Datei

% ── Metadaten ───────────────────────────────────────────────
\title{Titel der Präsentation}
\subtitle{Untertitel}
\author{Vorname Nachname}
\institute{Universität · Seminar: \textit{Seminarname}}
\date{\today}

% ============================================================
\begin{document}

% ── Titelfolie ──────────────────────────────────────────────
\maketitle

% ── Gliederung ──────────────────────────────────────────────
\begin{frame}{Gliederung}
  \tableofcontents
\end{frame}

% ============================================================
\section{Einführung}
% ============================================================

% Einfache Bullet-Liste
\begin{frame}{Thema und Motivation}
  \begin{itemize}
    \item Erster Punkt zur Einführung
    \item Zweiter Punkt – Kontext und Relevanz
    \item Ziel dieses Vortrags
  \end{itemize}
\end{frame}

% Block + Beispiel nebeneinander
\begin{frame}{Definition}
  \begin{block}{Kernbegriff}
    Hier steht eine formale oder inhaltliche Definition des zentralen Begriffs.
  \end{block}

  \bigskip
  \begin{example}
    Ein konkretes Beispiel, das die Definition veranschaulicht.
  \end{example}
\end{frame}


% ============================================================
\section{Lösungen}
% ============================================================

% TikZ-Ablaufdiagramm
\begin{frame}{Lösungsansatz – Überblick}
  \begin{center}
    \begin{tikzpicture}[
        node distance = 2cm,
        box/.style    = {draw, rounded corners, minimum width=2.8cm,
                         minimum height=0.7cm, align=center, font=\small},
        arr/.style    = {-Stealth, thick},
      ]
      \node[box] (a) {Schritt 1};
      \node[box] (b) [right=of a] {Schritt 2};
      \node[box] (c) [right=of b] {Schritt 3};
      \node[box] (d) [right=of c] {Ergebnis};
      \draw[arr] (a) -- (b);
      \draw[arr] (b) -- (c);
      \draw[arr] (c) -- (d);
    \end{tikzpicture}
  \end{center}
  \begin{itemize}
    \item Erster Schritt: \ldots
    \item Zweiter Schritt: \ldots
  \end{itemize}
\end{frame}

% Code-Folie (fragile ist Pflicht bei lstlisting!)
\begin{frame}[fragile]{Lösung – Codebeispiel}
  \begin{lstlisting}
-- Beispiel-Query
SELECT spalte1, spalte2
FROM   tabelle
WHERE  bedingung = TRUE
ORDER  BY spalte1;
  \end{lstlisting}
  \begin{itemize}
    \item Erläuterung Zeile 1--2: Was wird selektiert?
    \item Erläuterung Zeile 3: Filterbedingung
  \end{itemize}
\end{frame}

% ============================================================
\section{Kritik}
% ============================================================

% Alertblock für Schwächen
\begin{frame}{Limitierungen}
  \begin{alertblock}{Bekannte Schwächen}
    Dieser Ansatz funktioniert nicht optimal, wenn \ldots
  \end{alertblock}

  \bigskip
  \begin{itemize}
    \item Einschränkung 1 – Begründung
    \item Einschränkung 2 – Begründung
    \item Offene Frage / ungelöstes Teilproblem
  \end{itemize}
\end{frame}


% ============================================================
\section{Fazit}
% ============================================================

\begin{frame}{Zusammenfassung}
  \begin{itemize}
    \item \textbf{Problem:} Kurze Wiederholung der Fragestellung
    \item \textbf{Lösung:} Kernidee in einem Satz
    \item \textbf{Stärken:} Was funktioniert gut?
    \item \textbf{Schwächen:} Was bleibt offen?
  \end{itemize}

  \bigskip
  \begin{block}{Kernaussage}
    Der vorgestellte Ansatz eignet sich besonders für \ldots,
    stößt aber bei \ldots an seine Grenzen.
  \end{block}
\end{frame}

\begin{frame}{Ausblick}
  \begin{itemize}
    \item Mögliche Erweiterung 1
    \item Mögliche Erweiterung 2
    \item Interessante Folgefrage für künftige Arbeit
  \end{itemize}
\end{frame}

% ── Literatur ───────────────────────────────────────────────
\begin{frame}[allowframebreaks]{Literatur}
  \printbibliography[heading=none]
\end{frame}

\end{document}
